\section*{Results}
\subsection*{E1: detection skill with all labels}
All 20 training runs converged without numerical problems; the evaluation-selected checkpoint was reached after 504--1120 of the %%STEPS%% steps. Table~\ref{tab:main} lists test performance. Averaged over seeds, the CNN--LSTM obtained the highest ROC-AUC (%%AUC_CNNLSTM%%), followed by the CNN (%%AUC_CNN%%), the ConvLSTM (%%AUC_CONVLSTM%%) and the Transformer (%%AUC_TRANS%%). The Friedman test over seeds indicates that the ranking of the four models is not random ($\chi^2_3=%%FRIEDAUCCHI%%$, $p%%FRIEDAUCP%%$); for accuracy the evidence is weaker ($\chi^2_3=%%FRIEDACCCHI%%$, $p%%FRIEDACCP%%$).

The differences are, however, small compared with the sampling uncertainty of a %%NTEST%%-example test set. The seed-ensemble ROC-AUC values lie between %%EAUC_CONVLSTM%% and %%EAUC_CNNLSTM%%, and their bootstrap 95\,\% intervals overlap widely (Table~\ref{tab:main}, Fig.~\ref{fig:roc}). The largest pairwise difference in ensemble ROC-AUC is %%MAXDAUC%%, and no pairwise comparison is significant after Holm adjustment (%%MINPHOLM%% for both the bootstrap and the McNemar tests, and the smallest unadjusted bootstrap $p$-value is %%MINPRAW%%; Table~\ref{tab:pairs}). We therefore cannot claim that any architecture detects earthquakes better than another on clean data from this array.

Accuracy at the fixed 0.5 threshold was much less stable than ROC-AUC. Across seeds, CNN accuracy ranged from %%ACCMIN_CNN%% to %%ACCMAX_CNN%% while its ROC-AUC varied by only about $\pm0.01$; the ConvLSTM was the most stable (accuracy %%ACCMIN_CONVLSTM%%--%%ACCMAX_CONVLSTM%%). The low-accuracy runs ranked events correctly but placed many of them below 0.5, so the problem is the calibration of the output, not the ranking. With seed ensembling all four models reached 0.88--0.91 accuracy; the CNN--LSTM ensemble missed %%FN_CNNLSTM%% of 57 test earthquakes and produced %%FP_CNNLSTM%% false alarms in 57 noise windows.

For reference, the 2-D CNN of \citet{HuotArxiv2022}, trained on the full-resolution input (286 channels $\times$ 512 samples), reported 92.98\,\% accuracy and 0.9895 ROC-AUC on the test split. Our models, which see a six-fold spatially and two-fold temporally averaged input, are trained for 20 epochs on a CPU and contain fewer than $1.6\times10^5$ parameters, are below that reference. The present comparison is meant to compare architectures with each other under identical conditions, not to improve on the original detector.

\subsection*{E2: robustness to synthetic noise and dead channels}
Figure~\ref{fig:noise} and Table~\ref{tab:noise} show how the trained models respond to perturbations of the real test records. Additive white noise down to 5\,dB SNR changed ROC-AUC by less than 0.02 for every architecture. Below that the models separate. At 0\,dB the mean per-seed loss of ROC-AUC was %%N0DROP_CNN%% for the CNN, against %%N0DROP_CNNLSTM%%--%%N0DROP_TRANS%% for the other three models, although these differences were not significant (Friedman $p%%N0FRIEDP%%$). At $-5$\,dB the degradation differed clearly between architectures (Friedman $p%%N5FRIEDP%%$): the mean loss of ROC-AUC was %%N5DROP_CNN%% (CNN), %%N5DROP_CNNLSTM%% (CNN--LSTM), %%N5DROP_CONVLSTM%% (ConvLSTM) and %%N5DROP_TRANS%% (Transformer). The seed ensembles show the same order, with ROC-AUC of %%N5EAUC_CNN%%, %%N5EAUC_CNNLSTM%%, %%N5EAUC_CONVLSTM%% and %%N5EAUC_TRANS%%. The Transformer ensemble was significantly better than the CNN ensemble ($\Delta$AUC$_{\mathrm{CNN}-\mathrm{Transformer}}=%%N5D_CNN_TRANS%%$, 95\,\% CI %%N5CI_CNN_TRANS%%, Holm-adjusted $p%%N5P_CNN_TRANS%%$) and than the CNN--LSTM ensemble ($\Delta$AUC$_{\mathrm{CNN\text{-}LSTM}-\mathrm{Transformer}}=%%N5D_CNNLSTM_TRANS%%$, $p%%N5P_CNNLSTM_TRANS%%$), and the ConvLSTM ensemble was significantly better than the CNN ensemble ($\Delta$AUC$_{\mathrm{CNN}-\mathrm{ConvLSTM}}=%%N5D_CNN_CONVLSTM%%$, $p%%N5P_CNN_CONVLSTM%%$).

Removing channels was much less harmful. With 10\,\% or 30\,\% of the channels set to zero, ROC-AUC stayed within about 0.01 of the clean value for all models, and with 50\,\% of the channels removed the mean loss was only %%D50DROP_CONVLSTM%%--%%D50DROP_CNN%%, with no significant difference between architectures (Friedman $p%%D50FRIEDP%%$).

\subsection*{E3: label scarcity}
%%E3TEXT%%

\subsection*{Computational cost}
Table~\ref{tab:cost} summarises cost on one CPU thread. The Transformer had the most parameters (%%PAR_TRANS%%) but the fewest forward operations (%%GFLOP_TRANS%% GFLOPs per patch) and the shortest single-patch latency (%%LAT_TRANS%%\,ms). The ConvLSTM had the fewest parameters (%%PAR_CONVLSTM%%) but the highest operation count (%%GFLOP_CONVLSTM%% GFLOPs) and, because its 80 recurrent steps run in sequence, the longest latency (%%LAT_CONVLSTM%%\,ms) and training time (%%TIME_CONVLSTM%%\,s per run against about 325--340\,s for the other models). Peak memory during training with batch 32 increased by 285--380\,MB for all models. Even the slowest model processes one 25.6-s, 288-channel window in about 13\,ms on a single CPU thread, more than three orders of magnitude faster than real time.
